\IfFileExists{papers/formalized-vs-literature-preamble.tex}{\documentclass[11pt]{article}
\usepackage[utf8]{inputenc}
\usepackage[T1]{fontenc}
\usepackage{amsmath,amssymb,mathtools}
\usepackage{booktabs,tabularx,array}
\usepackage{enumitem}
\usepackage[margin=1in]{geometry}
\usepackage{xcolor}
\usepackage[hidelinks]{hyperref}
\usepackage{microtype}
\setlength{\parindent}{0pt}
\setlength{\parskip}{0.55em}
\newcommand{\Lean}[1]{\texttt{\nolinkurl{#1}}}
\newcommand{\RepoFile}[1]{\texttt{\nolinkurl{#1}}}
\newcommand{\Class}[1]{\textbf{#1}}
\newcommand{\Top}{\mathord{\top}}
\newcommand{\R}{\mathbb{R}}
\newcommand{\ENN}{\mathbb{R}_{\ge 0}^{\infty}}
\newcommand{\KL}{\operatorname{KL}}
\newcommand{\W}{\mathsf{W}}
\newcommand{\statusline}[2]{%
  \begin{center}\fbox{\begin{minipage}{0.92\linewidth}\textbf{Completion class:} #1\\[2pt]
  \textbf{Strongest caveat:} #2\end{minipage}}\end{center}}
\newenvironment{comparison}{\tabularx{\linewidth}{>{\bfseries}p{0.25\linewidth}X}}{\endtabularx}
\newcommand{\snapshot}{\texttt{902f51aa01a737af68f6feb71c83263cc9b1f4d9}}
}{\documentclass[11pt]{article}
\usepackage[utf8]{inputenc}
\usepackage[T1]{fontenc}
\usepackage{amsmath,amssymb,mathtools}
\usepackage{booktabs,tabularx,array}
\usepackage{enumitem}
\usepackage[margin=1in]{geometry}
\usepackage{xcolor}
\usepackage[hidelinks]{hyperref}
\usepackage{microtype}
\setlength{\parindent}{0pt}
\setlength{\parskip}{0.55em}
\newcommand{\Lean}[1]{\texttt{\nolinkurl{#1}}}
\newcommand{\RepoFile}[1]{\texttt{\nolinkurl{#1}}}
\newcommand{\Class}[1]{\textbf{#1}}
\newcommand{\Top}{\mathord{\top}}
\newcommand{\R}{\mathbb{R}}
\newcommand{\ENN}{\mathbb{R}_{\ge 0}^{\infty}}
\newcommand{\KL}{\operatorname{KL}}
\newcommand{\W}{\mathsf{W}}
\newcommand{\statusline}[2]{%
  \begin{center}\fbox{\begin{minipage}{0.92\linewidth}\textbf{Completion class:} #1\\[2pt]
  \textbf{Strongest caveat:} #2\end{minipage}}\end{center}}
\newenvironment{comparison}{\tabularx{\linewidth}{>{\bfseries}p{0.25\linewidth}X}}{\endtabularx}
\newcommand{\snapshot}{\texttt{902f51aa01a737af68f6feb71c83263cc9b1f4d9}}
}
\title{Finite Sinkhorn iteration convergence: formalized result versus the literature}
\author{}
\date{Formalization snapshot \snapshot}
\begin{document}
\maketitle
\statusline{\Class{S/R}: complete convergence theorem for explicitly packaged positive, gauge-normalized iterate systems}{The theorem starts from an iterate-system predicate and a named target potential system; it is not yet a single raw-initialization theorem with the exact Franklin--Lorenz sharp rate.}
\section{Literature targets}
The principal sources are Sinkhorn--Knopp (1967), Franklin--Lorenz (1989) Section~3, and the finite Fortet/IPF discussion in Chen--Georgiou--Pavon.  See
\RepoFile{prose/distilled_literature/SinkhornKnopp1967_matrix_scaling.tex} and
\RepoFile{prose/distilled_literature/FranklinLorenz1989_matrix_scaling.tex}.

\section{Lean capstone}
\Lean{ChenGeorgiouPavon2021.sinkhorn_iterates_converge_to_potentials} in
\RepoFile{ChenGeorgiouPavon2021/SocOt/Sinkhorn/Convergence/Gauge.lean} proves coordinatewise convergence of all four potential sequences.

Its inputs include:
\begin{itemize}[leftmargin=2em]
\item positive marginals and strictly positive finite kernel;
\item \Lean{IsFiniteSinkhornIterateSystem}, which packages positivity and the four update equations;
\item \Lean{IsFiniteSinkhornGaugeNormalized}, fixing the multiplicative gauge;
\item \Lean{IsFiniteSinkhornPotentialSystem}, identifying the target solution.
\end{itemize}

\section{Proof architecture}
The development proves uniform boxes, subsequence compactness, passage of mixed-phase equations to cluster points, uniqueness up to scaling, gauge-fixed uniqueness, and convergence from uniqueness of all cluster points.  A Franklin--Lorenz branch proves geometric decay of projective ratio spreads using the finite Birkhoff coefficient.  A separate projective-lag branch is retained as alternate compactness scaffolding.

\section{Comparison}
\begin{comparison}
Kernel & Strictly positive finite square kernel.\\
Raw initialization & Encoded indirectly by the iterate-system predicate rather than a constructor from one seed.\\
Existence of target & Available from the separate matrix-scaling theorem, but supplied to this capstone.\\
Gauge & Explicitly fixed, avoiding the false claim of convergence to an arbitrary representative.\\
Convergence & All four potential sequences converge.\\
Rate & Several geometric spread bounds are proved; the capstone does not package the exact source rate in its conclusion.\\
Zeros/support & Not handled.\\
\end{comparison}

\section{Next source-faithful strengthening}
Construct the iterate system from a raw positive seed, compose existence and uniqueness internally, expose convergence of the scaled matrices as well as potentials, and package the Franklin--Lorenz geometric constant and tail estimate in a theorem directly comparable with Theorem~4.
\end{document}
